\documentclass{article}

\usepackage[preprint]{neurips_2024}
\usepackage[utf8]{inputenc}
\usepackage[T1]{fontenc}
\usepackage{hyperref}
\usepackage{url}
\usepackage{booktabs}
\usepackage{amsfonts}
\usepackage{amsmath}
\usepackage{amssymb}
\usepackage{microtype}

\title{Algebraic Capability Tiers:\\
       The Cayley--Dickson Staircase of AI}

\author{Sounio Project}

\begin{document}
\maketitle

\begin{abstract}
Classical ML sits comfortably on the real line.  Complex, quaternion,
and octonion generalizations have been proposed for signal
processing, 3D rotation learning, and associative-alternative
networks.  In Papers~B--H we showed that the \emph{sedenion} level
(Cayley--Dickson level~4) is the \emph{first} algebra on the tower
in which zero divisors exist, and argued that this single structural
fact enables surgical machine unlearning, model editing, and
capability gating as algebraic identities rather than empirical
tricks.  Here we take one step further up the tower.  We define the
\textbf{Cayley--Dickson Capability Ladder}: a 5-tier taxonomy of
algebraic structure---$\mathbb R$, $\mathbb C$, $\mathbb H$,
$\mathbb S$, $\mathbb T$ (pathions, 32-dim)---and exhibit, for
tier~L5, a strict extension of the sedenion zero-divisor graph:
the primitive two-support basis is $5\times$ larger (420 vs.\ 84),
and the Fano-line analogue (xor-fiber label count) grows from
7~(PG(2,2)) to 15~(PG(3,2)).  Both counts are verified by Lean~4
\texttt{native\_decide} on the committed enumeration
(\texttt{SounioPathionBridge.lean}).  We interpret L5 as the
first tier on which \emph{meta-surgical} operations become
algebraically available: surgery on the surgery operator itself.
No claim is made that any existing model \emph{uses} level~5; the
contribution is the algebraic inventory, the tier taxonomy, and the
formal verification of the L4$\to$L5 growth rates.
\end{abstract}

\section{Introduction}

The Cayley--Dickson doubling construction produces an infinite
tower of normed algebras

\begin{center}
  \begin{tabular}{ccccccc}
  $\mathbb R$ &$\to$& $\mathbb C$ &$\to$& $\mathbb H$ &$\to$& $\mathbb O$ \\
   L0 & & L1 & & L2 & & L3
  \end{tabular}
  \vspace{-0.15em}
  \begin{tabular}{ccccccc}
  $\to$& $\mathbb S$ &$\to$& $\mathbb T$ &$\to$& \ldots \\
       & L4 & & L5
  \end{tabular}
\end{center}

Each step doubles the dimension.  Each step also *loses* a property:
L1 loses ordering, L2 loses commutativity, L3 loses associativity,
L4 loses alternativity \emph{and} gains zero divisors.  Paper~B
identified the L3$\to$L4 threshold as the \emph{first} point at
which the ZD structure enables exact machine unlearning.  Papers
C--G extended the surgical calculus to six operations.  Paper~H
extended it to training dynamics.  This paper asks: what happens at
L5?

\paragraph{Methodology disclosure.}
We tag every structural claim in this paper with one of two markers.
\textbf{[LEAN]} --- mechanically verified by \texttt{lake build} against
\texttt{SounioZeroDivisorBridge.lean} (L4: 84 primitives, 168 classes,
7 fibres) and \texttt{SounioPathionBridge.lean} (L5: 420 primitives,
15 fibre labels); every count is closed by \texttt{native\_decide} over
concrete Cayley--Dickson multiplication, with no \texttt{sorry}, no
Mathlib, and no \texttt{omega}-fall-through.
\textbf{[STRUCTURAL]} --- algebraic facts (loss of ordering, commutativity,
associativity, alternativity) that follow from the Cayley--Dickson
doubling construction and Hurwitz's theorem; they are not measured but
are well-established in the literature we cite.
This paper makes no claim about an L5 (pathion) ZD class count or L5-based
operational benchmark; those are explicit future work.

\section{The ladder}

\begin{table}[h]
  \centering
  \caption{[LEAN] + [STRUCTURAL] The Cayley--Dickson Capability Ladder.  ``Surgical'' = the
           algebra admits the six Sounio surgical types; ``Meta'' =
           the algebra admits surgery on the surgical operator
           itself.}
  \label{tbl:ladder}
  \begin{tabular}{cllcccc}
    \toprule
    Tier & Algebra      & Dim  & Assoc & Alt & ZD    & Surgical / Meta \\
    \midrule
    L0   & $\mathbb R$   & 1   & yes   & yes & no    & --- / ---        \\
    L1   & $\mathbb C$   & 2   & yes   & yes & no    & --- / ---        \\
    L2   & $\mathbb H$   & 4   & yes   & yes & no    & --- / ---        \\
    L3   & $\mathbb O$   & 8   & no    & yes & no    & --- / ---        \\
    L4   & $\mathbb S$   & 16  & no    & no  & yes   & surgical / ---   \\
    L5   & $\mathbb T$   & 32  & no    & no  & yes   & surgical / meta  \\
    \bottomrule
  \end{tabular}
\end{table}

The stdlib exposes the tier view via
\texttt{stdlib/algebra/ladder.sio}:
\texttt{tier\_of\_bits(k)} returns a struct with fields
\texttt{assoc}, \texttt{alternative}, \texttt{has\_zd},
\texttt{surgical}, \texttt{meta\_surgical}.
Downstream code can branch on the tier (e.g.\ a training pipeline
refuses to instantiate an \texttt{ExactlyPrivate<T>} head on an L3
algebra because \texttt{has\_zd} is false, producing a clean
compile-time error message rather than silent failure).

\section{Level-5 primitive census}

At level~5, the pathion imaginary basis has 31 elements
($e_1,\ldots,e_{31}$).  A two-support primitive element is
$e_\ell + s\,e_h$ with
\begin{itemize}
  \item $\ell \in \{1, \ldots, 15\}$ (imaginary sedenion basis),
  \item $h \in \{17, \ldots, 31\}$ (pathion extension),
  \item $\ell \oplus h \neq 16$ (excludes the diagonal family),
  \item $s \in \{+1, -1\}$.
\end{itemize}

\textbf{Census theorem (Lean, \texttt{native\_decide}).}  The number
of valid primitive pathion two-support elements is exactly
$$
|\text{pathPrims}| \;=\; 420.
$$

\textbf{Fiber-label theorem (Lean, \texttt{native\_decide}).} The
xor-fiber label set has cardinality
$$
|\text{pathFiberLabels}| \;=\; 15,
$$
every label lying in $\{17,\ldots,31\}$.  This matches the line
count of PG(3,2), the 3-dim projective space over $\mathbb F_2$,
extending the PG(2,2) = Fano plane structure of level~4.

\textbf{Ratio theorem.}  The primitive count at L5 is exactly
$5\times$ the L4 count, and the fiber count strictly grows:
$420 = 5\times 84$ and $15 > 7$.

All three theorems live in
\texttt{formal/lean4/SounioPathionBridge.lean}.  No \texttt{sorry},
no \texttt{Mathlib}, all counts closed by \texttt{native\_decide}
on committed finite enumerations.

\begin{table}[h]
  \centering
  \caption{[LEAN] Two quantities that strictly grow from L4 to L5.}
  \label{tbl:growth}
  \begin{tabular}{lcc}
    \toprule
    Quantity                   & Sedenion (L4) & Pathion (L5) \\
    \midrule
    Imaginary-basis dim        & 15            & 31           \\
    Primitive two-support count & 84            & 420          \\
    Xor-fiber label count       & 7             & 15           \\
    Matching projective space  & PG(2,2)       & PG(3,2)      \\
    ZD class count (projective) & 168           & bounded$^*$  \\
    \bottomrule
  \end{tabular}\\
  \footnotesize
  $^*$ A tight \texttt{native\_decide} count for the pathion ZD
  class cardinality is left as a CI-scope computation; only an
  upper bound is checked at library-build time.
\end{table}

\section{Why tiers matter for AI}

The Sounio type-system story is that algebraic \emph{tiers} fence
off capabilities that lower tiers cannot express.  At L4 the new
capability is \emph{surgery} (exact unlearning, etc).  At L5 the
new capability, we argue, is \emph{meta-surgery}: surgery on the
surgical operator.  In the pathion structure, a primitive
two-support element lives on a \emph{line} of fibers inherited from
PG(3,2), i.e.\ one sedenion surgery-operator identity $A \in \mathbb S$
naturally belongs to a whole \emph{1-parameter family} of pathion
surgery-identities obtained by lifting the L4 label along the L5
fiber.  That family is a Lean-enumerable list; navigating it is
literally ``choosing which surgery to do''.

We do not commit here to a specific downstream model that exploits
L5; the point is that the algebraic inventory and the tier taxonomy
now exist, formally verified, and are available to downstream
researchers through the \texttt{stdlib/algebra/ladder.sio} API.  An
example of how Paper-I structure might be used: a future compiler
lint rule could forbid \texttt{meta\_surgical} operations inside
standard L4 modules, producing a compile-time error analogous to
the \texttt{with ZD} check of Paper~C.

\section{Related algebraic work}

Previous work on higher-order hypercomplex neural networks is
largely restricted to L1 (complex-valued CNNs~\cite{trabelsi2018deep})
and L2 (quaternion networks for 3D~\cite{parcollet2020quaternion}).
Octonion networks (L3) have been explored in signal
processing~\cite{octocomp2022}.  Sedenion networks (L4) have been
proposed mainly for associative memory~\cite{sedenionmem2021}, not
for surgical guarantees.  We know of no prior neural-network or
symbolic-AI work that targets L5.  The companion compiler work is
at~\cite{sounio2026zd}.

\section{What this paper does not claim}

\begin{itemize}
  \item No empirical claim about pathion-based neural networks.
  \item No tight ZD-class count for L5 (only a bound; tight count
        is a CI-scope computation that requires cheaper ZD-pair
        enumeration strategies).
  \item No claim that L6 tiers are obviously useful; the paper
        stops at L5 because that is where meta-surgery first
        becomes definable.
  \item No claim that the dimensional cost ($\times 2$ per tier) is
        ever justified in a deployed model; the practical sweet
        spot may well remain L4.
\end{itemize}

\section{Conclusion}

We gave a formal, machine-checked extension of the Sounio
zero-divisor bridge from the sedenion level (L4, 16-dim) to the
pathion level (L5, 32-dim).  At L5 the primitive two-support basis
grows to $420 = 5\times 84$, and the xor-fiber structure lifts from
Fano PG(2,2) to PG(3,2) with 15 lines.  Both counts are closed by
\texttt{native\_decide} in Lean~4.  The \texttt{stdlib/algebra/ladder.sio}
module exposes a 5-tier taxonomy that downstream Sounio code can
query.  The Cayley--Dickson ladder gives Sounio a principled way
to talk about algebraic \emph{capability tiers} for AI.  Paper~I
completes the theoretical part of the Butterfly Thesis; the
practical artifact (ZD-SSM) and the self-editing easter egg follow
in Papers J and K.

\bibliographystyle{plain}
\begin{thebibliography}{9}
\bibitem{trabelsi2018deep}
C. Trabelsi et al.
\emph{Deep complex networks.}
ICLR, 2018.

\bibitem{parcollet2020quaternion}
T. Parcollet, M. Morchid, and G. Linarès.
\emph{A survey of quaternion neural networks.}
Artificial Intelligence Review, 2020.

\bibitem{octocomp2022}
J. Ward.
\emph{Octonion-valued deep learning: a survey.}
Neural Processing Letters, 2022.

\bibitem{sedenionmem2021}
E. Saoud et al.
\emph{Sedenion-based neural associative memories.}
Neurocomputing, 2021.

\bibitem{sounio2026zd}
Sounio Project.
\emph{Paper B: the sedenion zero-divisor structure of AI surgery.}
Companion preprint, 2026.
\end{thebibliography}

\end{document}
